\begin{textsample}{POS Dim 1 – es – Score 44.00 – es/es\_em\_14.txt}  \label{ex:f1_pos_133}
A ÁREA DE LINGUAGENS E SUAS TECNOLOGIAS “ A linguagem e a vida são \textbf{uma} coisa só [... ] ; e como a vida é \textbf{uma} corrente contínua, a linguagem também deve evoluir constantemente ” – Guimarães Rosa exprime, em entrevista a Günter Lorenz, o quão os seres humanos estão submersos nos mais diversos “ mundos ” da “ contínua linguagem ”.
A linguagem é \textbf{um} dos aspectos do fenômeno da comunicação humana, sendo considerada \textbf{um} conceito amplo \textbf{que} engloba quaisquer códigos utilizados pelos \textbf{sujeitos}, a fim de interagirem \textbf{uns} com os outros e, por meio dela, pode ser estabelecida comunicação entre dois ou mais interlocutores.
Embora não seja privilégio apenas dos seres humanos, evidencia as multifacetadas particularidades \textbf{que} eles têm não só no \textbf{que} tange à sua interação com os da mesma espécie, mas – e também – no \textbf{que} se refere à sua sobrevivência.
Da mesma maneira como \textbf{que} ser “ racional ” diferencia os \textbf{homens} dos outros animais, a linguagem como fenômeno também o faz, porém de formas mais complexas e, ao mesmo tempo, mais \textbf{inerentes} às particularidades humanas.
A linguagem humana é a expressão do pensamento.
Materializamos nossas relações por meio da linguagem.
Ela cria oportunidades de convívio e interação social entre indivíduos, por meio das modalidades oral, escrita, gesticulada, imagética, dispostas em diversos gêneros.
Assim como \textbf{salienta} Guimarães Rosa, a linguagem é vida – vida constituída por símbolos, sinais, regras, convenções e situações de interação \textbf{que} permitem a \textbf{inter-relação} entre os \textbf{homens} : a expressão de ideias, sentimentos, sensações, bem como a intervenção do indivíduo no mundo e o registro histórico dessas manifestações.
Além disso, por meio da linguagem podemos associar sinais, tais como palavras, desenhos, paisagens, gestos e textos em múltiplos suportes, com as ideias e com os conceitos \textbf{que} se formam em nosso cérebro.
Sendo assim, a linguagem possui função de sua importância no percurso histórico da \textbf{humanidade}.
A Área de Linguagens e suas Tecnologias, instrumento e meio de registro histórico da \textbf{humanidade}, deve abarcar os nove princípios do Artigo 5º, da Resolução No 3, de 21 de novembro de 2018, do Conselho Nacional de Educação, da Câmara de Educação Básica ( CNE/CEB ), dentre os quais destacamos : formação integral do estudante, expressa por valores, aspectos físicos, cognitivos e socioemocionais, para \textbf{que} sejam capazes de se colocar no lugar do outro, ser colaborativos, flexíveis e sensíveis às questões coletivas.
Dessa forma, o desenvolvimento socioemocional não pode constar em \textbf{uma} parte isolada dos textos do currículo, nem pode aparecer neles de forma subentendida, pouco valorizada ou apenas a serviço da aprendizagem de conteúdos.
Os estudos nessa área devem, portanto, visar ao aprofundamento de conhecimentos estruturantes para aplicação de diferentes linguagens em contextos sociais e de trabalho, estruturando arranjos curriculares \textbf{que} permitam estudos nas línguas maternas ( tanto o português como a língua pomerana, os dialetos quilombolas e falares indígenas – considerando, \textbf{aqui}, as diversidades culturais e regionais do Espírito Santo ) nas línguas estrangeiras, clássicas, na Língua Brasileira de Sinais ( LIBRAS ), das artes, design, linguagens digitais, corporeidade, artes cênicas, roteiros, produções \textbf{literárias}, dentre outras.
No Ensino Médio, a compreensão das linguagens está vinculada às práticas sociais e aos discursos \textbf{que} circulam socialmente.
Nesse sentido, o reconhecimento da pluralidade linguística \textbf{que} \textbf{permeia} as práticas sociais deverá estimular o respeito, o reconhecimento da democracia e da igualdade, por meio da compreensão dos processos identitários intrínsecos às experiências concretas de linguagem.
Dessa mesma maneira, os processos comunicacionais advindos dos usos de mundo e dos usos da linguagem \textbf{aqui} pensados têm por intuito levar em consideração as multifacetadas práticas sociais, possibilitando aos indivíduos o exercício da autonomia plena de suas vivências, de suas experiências, dos seus vínculos de afetividade e das suas trocas.
Acreditamos, assim, em \textbf{um} processo comunicativo \textbf{baseado} em relações fecundas.
O vocábulo comunicação é proveniente do latim communicare, \textbf{que} significa “ tornar comum ”, “ compartilhar ”.
O conceito de pluralidade linguística \textbf{aqui} veiculado visa a permitir \textbf{que} a linguagem funcione como \textbf{uma} espécie de mola propulsora de participação do estudante nos meandros \textbf{que} circundam a sua juventude, as suas particularidades, as suas experiências e, principalmente, a sua autonomia.
É sabido \textbf{que} a juventude, \textbf{enquanto} construção sociocultural, é marcada pela diversidade e pelo multiculturalismo.
Quando nos referimos aos jovens estudantes do Ensino Médio do Espírito Santo, não podemos desconsiderar as várias identidades juvenis \textbf{que} emergem nessa sociedade, \textbf{ora} subvertendo as regras e padrões impostos, \textbf{ora} \textbf{exigindo} aceitação, respeito, ajuda e solidariedade.
Sendo assim, não podemos falar em apenas \textbf{uma} juventude, mas em “ juventudes ”, \textbf{uma} vez \textbf{que} esses jovens, embora em idade similar, são heterogêneos e \textbf{vivem} diferentes realidades e contextos no cotidiano capixaba.
A juventude, por ser \textbf{uma} fase de incertezas e de construção de identidades, o preconceito ( de etnia, de gênero, de classe, entre outros ) é a \textbf{marca} mais profunda na vida de muitos jovens.
A falta de aceitação das diversas culturas, a desigualdade social, a \textbf{exclusão} e a violência potencializam a discriminação, de modo \textbf{que} muitos ficam à margem da sociedade e subjugados a expectativas negativas quanto ao futuro, tampouco possuem espaços e momentos para reflitirem sobre seu projeto de vida.
Nesse sentido, cabe à escola – \textbf{talvez} a única instituição \textbf{que} possa entender e acolher esse jovem na sua \textbf{singularidade} – permitir o diálogo entre as várias nuances desse \textbf{sujeito} em formação, bem como possibilitar o desenvolvimento integral do estudante, oferecendo tais espaços e momentos.
As vozes das juventudes capixaba ecoam de forma múltipla e polifônica, oriundas do campo e da cidade, do mar e das montanhas.
É necessário \textbf{que} a escola permita \textbf{que} essa multiplicidade de vozes se instaure de forma igualitária.
Por esse motivo, a Área de Linguagens e Suas Tecnologias tem como objetivo possibilitar o contato e a valorização das diversas manifestações culturais e artísticas, ampliando o repertório cultural dos envolvidos no processo de ensino/aprendizagem e permitindo reflexões sobre as linguagens artísticas, verbais e corporais.
No componente curricular Língua Portuguesa, priorizamos a competência leitora, cuja parte integrante mais significativa é o texto ; afinal, é partindo das mais diversas tipologias textuais \textbf{que} são estabelecidos os estudos dos eixos \textbf{norteadores} dos campos de saber \textbf{que} envolvem a nossa língua materna : tanto gramaticais e linguísticos, quanto interpretativos e \textbf{literários}.
Em outros termos, tendo -se em vista as demandas educacionais em voga no presente momento, \textbf{salientamos} a necessidade de \textbf{um} trabalho sistemático com a oralidade nos mais diversos contextos comunicativos, com a (re)escrita, com a (re)leitura e com a experimentação de práticas culturais, artísticas, \textbf{literárias}, midiáticas e tecnológicas relevantes \textbf{que} envolvam as variedades linguísticas e as significativas multiplicidades de gêneros discursivos, bem como os multiletramentos e os novos letramentos.
Espera -se, nos parâmetros da \textbf{atual} conjuntura, \textbf{que} o profissional de Língua Portuguesa não só estabeleça diálogos interdisciplinares, com os componentes de sua área de conhecimentos, mas – e principalmente – diálogos transdisciplinares e multidisciplinares, \textbf{focados} nos temas integradores.
Nesse sentido, é indispensável, também, \textbf{que} os estudantes se sintam familiarizados e socialmente engajados nos usos e estudos da língua, levando -se em consideração os aspectos locais, culturais, comportamentais, artísticos e sociais específicos da sua região.
Esperamos \textbf{que}, ao fim da jornada do Ensino Médio, o componente curricular Língua Portuguesa estimule o jovem capixaba a exercer \textbf{uma} participação ativa e ética nas múltiplas práticas linguísticas e socioculturais.
O componente curricular Língua Espanhola, como língua estrangeira no Ensino Médio, possibilita a expansão dos repertórios linguísticos e culturais dos estudantes, bem como o acesso ao conhecimento científico e tecnológico produzido nessa língua.
Do mesmo modo, amplia seu \textbf{horizonte} de oportunidades no \textbf{que} se refere às possibilidades de \textbf{um} aprendizado contínuo, no campo da interpretação crítica da realidade.
Seu estudo, como língua estrangeira, possibilita a formação de \textbf{um} cidadão capaz de pensar o mundo com criticidade e conviver harmonicamente com a diversidade e o trânsito intercultural no qual está inserido, por ser a Língua Espanhola o idioma oficial de quase \textbf{totalidade} dos países \textbf{que} compõem a América do Sul.
As aprendizagens em Espanhol permitem \textbf{que} os estudantes usem esses conhecimentos com o objetivo de aprofundar a compreensão sobre o mundo em \textbf{que} \textbf{vivem}, bem como explorar novas perspectivas de pesquisa e obtenção de informações ; auxilia na administração de opiniões divergentes, de forma crítica, mas com respeito ao outro.
O amplo e vasto universo cultural e social \textbf{que} \textbf{permeia} os conhecimentos desse idioma possibilita aos jovens compreender, reconhecer e vivenciar valores e identidades e, por \textbf{conseguinte}, conviver democraticamente, respeitando toda essa diversidade.
A rede de ensino do Estado do Espírito Santo, considerando o fato de \textbf{que} já oferta Língua Espanhola conforme legislação vigente e por ter \textbf{um} número considerável de professores habilitados para trabalhar com esse componente curricular decidiu acatar o \textbf{que} é instruído no inciso IX, da Resolução No 3, de 21 de dezembro de 2018, do Conselho Nacional de Educação e oferecer a possibilidade de acesso aos conhecimentos dessa língua na parte flexível de seu currículo.
Por fim, estudar Espanhol possibilita aos jovens apreciarem a multiplicidade da produção artística e cultural, a partir da aprendizagem de suas múltiplas linguagens, assim como promove o desenvolvimento de sua competência produtiva, a partir da prática sistemática da língua.
O componente curricular Língua Inglesa ( LI ), de forma contextualizada, vem ao encontro do status de língua de comunicação entre os povos.
Essa prática aparece no Ensino Médio como \textbf{um} aprofundamento das habilidades já desenvolvidas no Ensino Fundamental.
É de \textbf{suma} importância \textbf{que} o ensino de Língua Inglesa considere o contexto do estudante, \textbf{uma} vez \textbf{que} a desmistificação das crenças relacionadas ao ensino da Língua Inglesa é necessária para acabar de vez com a ideia de \textbf{que}, para se tornar \textbf{um} bom falante da língua, é \textbf{imprescindível} falar como \textbf{um} nativo de países de LI.
Ao levar isso em consideração, o professor deve \textbf{focar} sua metodologia no ensino de \textbf{um} mundo híbrido, multicultural e digitalmente conectado no qual \textbf{vivemos}.
Deve perceber, e mostrar a seus alunos, \textbf{que} respeitando os aspectos locais dos \textbf{aprendizes} e suas identidades, há a promoção da inclusão social e, ao proporcionar o acesso e a produção de informações por meio da Língua Inglesa, consequentemente contribui com a formação de \textbf{um} cidadão global, multicultural e multiletrado, no qual o estudante percebe \textbf{que} os multiletramentos podem ser compreendidos como práticas sociais de linguagem realizadas em diversos contextos.
Ao alinhar o ensino de LI com a Base Nacional Comum Curricular – BNCC ( 2017 ), na qual é reestabelecido o \textbf{compromisso} com a formação integral do estudante, já previsto em marcos legais anteriores, o currículo de Língua Inglesa do Espírito Santo será \textbf{um} facilitador para o entendimento de \textbf{uma} proposta de educação linguística consciente e crítica, com foco na função social e política do inglês e, portanto, essa língua assume as concepções de língua franca e de língua adicional.
Vale mencionar \textbf{que} língua franca é o uso da língua inglesa entre falantes de diferentes línguas em contextos multilíngues, e, nesses contextos, a escolha de qual variedade da língua inglesa a ser usada \textbf{depende} dos falantes envolvidos.
Já a língua adicional é a língua estrangeira acrescida ao aprendizado desse aluno, além da sua língua materna.
O Espírito Santo tem \textbf{uma} função sócio, histórico e cultural muito importante, no \textbf{que} se refere ao ensino de LI, \textbf{uma} vez \textbf{que} é \textbf{um} estado \textbf{que} possui vocações turísticas e de comércio exterior, motivo pelo qual o desenvolvimento de habilidades e competências, nesse componente, possibilita ao estudante expandir suas possibilidades comunicativas e sociais.
O componente curricular Arte contribui de forma significativa para o desenvolvimento do protagonismo, da autonomia reflexiva, criativa e \textbf{expressiva} dos estudantes.
Considerando sua amplitude, promove o conhecimento do sujeito sobre si mesmo e sobre o outro, presentes em \textbf{um} mesmo cenário e/ou em seu entorno.
Nesse processo de aprendizagem, pesquisa e fazer artístico as \textbf{percepções} e compreensões do mundo se ampliam e se interconectam, numa perspectiva crítica, sensível e poética em relação à vida humana, permitindo aos \textbf{sujeitos} maior receptividade às \textbf{percepções} e experiências, \textbf{inspirando} a capacidade de imaginar e ressignificar o seu dia a dia, valorando os saberes dos estudantes.
O Currículo do Espírito Santo, considerando a BNCC, compreende as aprendizagens no Ensino Médio, como proposta de progressão, análise e aprofundamento na pesquisa e no desenvolvimento de processos de criação autorais nas linguagens : artes visuais, audiovisual, dança, teatro, artes circenses e música, visando a promover o engajamento dos estudantes nos processos criativos, oportunizando conexões e incorporação de estudos, pesquisas e referências estéticas, poéticas, sociais, culturais e políticas – nos âmbitos local, regional, nacional e internacional – para a criação de projetos artísticos individuais, coletivos e colaborativos.
O trabalho com a Arte no Ensino Médio convida ao entrelaçamento de culturas e saberes das distintas manifestações culturais \textbf{populares}, utilizando a linguagem e suas tecnologias de forma integrada, propiciando aos estudantes conhecimento, apropriação e valorização do patrimônio cultural, possibilitando análise crítica e problematizadora, estabelecendo relações entre arte, mídia, política, mercado e consumo, \textbf{aprimorando} sua capacidade de elaboração de análises em relação às produções estéticas \textbf{que} observam/vivenciam e criam e, consequentemente, promovendo maior respeito para com as diversas culturas \textbf{constituintes} da sociedade capixaba, dentre elas a cultura do campo, a cultura indígena e a cultura quilombola e, por extensão, a brasileira, motivando o protagonismo dos estudantes, como apreciadores, críticos e artistas, desenvolvendo trabalhos de criação e curadoria, de modo consciente, ético, e autônomo, nas diversas linguagens, em manifestações e/ou eventos artísticos e culturais, utilizando materiais, instrumentos e recursos convencionais, alternativos e digitais, em diferentes meios e tecnologias.
Finalmente, o componente curricular Educação Física, por meio da perspectiva da cultura corporal de movimento, garante sua contribuição na formação do \textbf{sujeito} e na construção da postura reflexiva diante do mundo, \textbf{uma} vez \textbf{que} permite explorar o movimento e a gestualidade em práticas corporais de diferentes grupos culturais e analisar os discursos e valores associados a elas.
A proposta de progressão das aprendizagens no Ensino Médio prevê o aprofundamento por \textbf{intermédio} dos 6 ( seis ) objetos de conhecimentos já abordados ao longo do Ensino Fundamental.
Desse modo, nessa etapa, os estudantes devem ser desafiados a refletir sobre jogos e brincadeiras, esportes, danças, lutas, ginásticas e práticas corporais de aventura – \textbf{aqui} denominados Práticas Corporais –, aprofundando seus conhecimentos, potencialidades e os limites do corpo, a importância de se assumir \textbf{um} estilo de vida ativo e seus impactos na manutenção da saúde.
Além disso, na Educação Física do Ensino Médio, esses fatos justificam tanto o uso das tecnologias digitais da informação e comunicação – TDIC, atreladas aos campos de atuação e, ao mesmo tempo, às manifestações da identidade social, ao bem-estar e à longevidade. \textbf{Logo}, é necessário assegurar às juventudes o pertencimento às práticas sociais e culturais locais, o exercício da autonomia e da cidadania, sem desconsiderar a cultura digital e os multiletramentos valorizados pela sociedade.
Para tal, é \textbf{imprescindível} \textbf{que} os cinco campos de atuação social, sejam priorizados na Área de Linguagens e suas Tecnologias, considerando -se \textbf{que} : O campo da vida pessoal possibilita ao jovem o reconhecimento pleno da sua autonomia na construção da sua identidade individual e social, resgatando as suas trajetórias e memórias, o seu autoconhecimento e alteridade para a estruturação de seus projetos de vida.
O campo das práticas de estudo e pesquisa almeja \textbf{uma} juventude preparada para lidar com o saber acadêmico-científico por meio da pesquisa, criação e construção de novos conhecimentos e, consequentemente, suscitada a aprender a aprender.
O campo jornalístico-midiático oportuniza \textbf{que} nossos jovens desenvolvam não só afinidades e familiaridades com os meios jornalísticos e midiáticos, como também lhes seja aguçada a consciência crítica perante a sociedade.
No campo de atuação na vida pública, a participação social torna -se evidenciada com o \textbf{estímulo} à condução do estudante a \textbf{uma} convivência ética e respeitosa entre os cidadãos e com a apropriação de gêneros legais e jurídicos/normativos utilizados em prol da defesa dos direitos do indivíduo, garantido o protagonismo em \textbf{face} da sua comunidade.
O campo artístico ( considerado em Língua Portuguesa como campo artístico-literário ) propicia ao jovem do Ensino Médio a ampliação/valorização da sensibilidade, da fruição estética e das experiências de processos criativos na construção de sua identidade e no (re)conhecimento da diversidade cultural e linguística \textbf{que} o circunda.
Enfim, o \textbf{que} se pretende com a organização desses campos é contribuir para a formação de jovens autônomos, participativos e éticos, de forma associada ao \textbf{compromisso} com a constituição de \textbf{uma} sociedade equitativa, e, assim como propõe Guimarães Rosa, em trecho já mencionado, ansiamos por jovens \textbf{que} entendam e associem a multiplicidade de linguagens à vida, buscando \textbf{uma} corrente contínua de evolução individual e cidadã.
A área de linguagens é composta por sete competências específicas \textbf{que} estão em articulação com as competências gerais da Educação Básica e com as da área de Linguagens do Ensino Fundamental apresentadas a seguir :

% matched lemmas: aprendiz, aprimorar, aqui, atual, basear, compromisso, conseguinte, constituinte, depender, enquanto, estímulo, exclusão, exigir, expressivo, face, focar, homem, horizonte, humanidade, imprescindível, inerente, inspirar, inter-relação, intermédio, literário, logo, marca, norteador, ora, percepção, permear, popular, que, salientar, singularidade, sujeito, sumo, talvez, totalidade, um, vivar, viver
\end{textsample}
